%==============================================================================
%  CS 601 -- AI and Uncertainty Reasoning
%  Per-topic POST-study assessment (3 questions per topic).
%
%  Parallel form to the pre-study assessment: identical topic list, identical
%  per-topic structure (Q1 single-answer, Q2 select-all, Q3 two-part
%  calculation) and identical point weights, so pre/post scores are directly
%  comparable for learning-gain analysis.  All questions and numbers are new.
%
%  Item-writing corrections carried over from the pre-study review:
%    * correct-answer positions are spread across A-D (not clustered on A);
%    * the correct option is not systematically the longest;
%    * the single false option in each "select all" item is scattered in
%      position and is not systematically the shortest;
%    * notation is glossed, acronyms are spelled out, graph load is minimised.
%
%  SINGLE SOURCE OF TRUTH.  The same file builds BOTH copies; the only switch
%  is the exam-class "answers" option:
%     Teacher's copy : pass "answers"  -> correct answers + worked solutions
%     Student's copy : omit it         -> questions only (all MCQ)
%
%  Build (run each twice so "Page x of N" resolves):
%     pdflatex -jobname=post-teacher "\PassOptionsToClass{answers}{exam}\input{post.tex}"
%     pdflatex -jobname=post-student "\input{post.tex}"
%==============================================================================

\documentclass[answers,addpoints]{exam}

\usepackage[utf8]{inputenc}
\usepackage[T1]{fontenc}
\usepackage{amsmath,amssymb}
\usepackage[margin=1in]{geometry}

% ---- Layout / headers --------------------------------------------------------
\pagestyle{headandfoot}
\firstpageheader{CS 601 -- AI and Uncertainty Reasoning}{}{Post-study Assessment}
\runningheader{CS 601 -- Post-study Assessment}{}{%
  \ifprintanswers TEACHER'S COPY\else STUDENT'S COPY\fi}
\firstpagefooter{}{Page \thepage\ of \numpages}{}
\runningfooter{}{Page \thepage\ of \numpages}{}

\checkboxchar{$\Box$}
\CorrectChoiceEmphasis{\bfseries}
\SolutionEmphasis{\itshape}

\newcommand{\topic}[2]{%
  \vspace{1.0\baselineskip}%
  \noindent\textbf{\large Topic #1: #2}\par
  \noindent\rule{\linewidth}{0.4pt}\par
  \vspace{0.4\baselineskip}}

\begin{document}

% ---- Cover block -------------------------------------------------------------
\begin{center}
  {\LARGE\bfseries CS 601 -- AI and Uncertainty Reasoning}\\[4pt]
  {\Large Per-topic Post-study Assessment}\\[6pt]
  {\large \ifprintanswers \textbf{TEACHER'S COPY}\\(answers, accepted ranges, and worked solutions)%
          \else \textbf{STUDENT'S COPY}\\(select the best option for every part)\fi}
\end{center}

\vspace{0.5\baselineskip}
\noindent
\textbf{Instructions.} Every part is multiple choice -- just select an option;
\emph{no written working is needed}. Each topic has three questions:
\textbf{Q1} single-answer multiple choice (choose ONE),
\textbf{Q2} ``select all that apply'' (choose every correct option), and
\textbf{Q3} a calculation in two parts: part~\textbf{(a)} asks \emph{which
expression is correct}, and part~\textbf{(b)} selects the value that expression
gives. Each question is self-contained. This post-study paper mirrors the
pre-study one in length and format.

\ifprintanswers
\noindent\textbf{Grading note (teacher's copy).} Each calculation part has a
single correct option; the intended exact value is shown in the worked solution.
\fi

\vspace{0.5\baselineskip}

%==============================================================================
\topic{1}{Probabilistic Reasoning}
%==============================================================================

\begin{questions}

% ---- Q1 (correct: C) ----
\question[2] (Single answer) For two events $A$ and $B$ with $P(B)>0$, which
equation correctly states \textbf{Bayes' rule}?
\begin{choices}
  \choice $P(A\mid B)=P(A)\,P(B)$
  \choice $P(A\mid B)=P(A,B)\,P(B)$
  \CorrectChoice $P(A\mid B)=\dfrac{P(B\mid A)\,P(A)}{P(B)}$
  \choice $P(A\mid B)=\dfrac{P(B\mid A)}{P(A)}$
\end{choices}
\begin{solution}\textbf{C.} Bayes' rule inverts a conditional: $P(A\mid B)=P(B\mid A)P(A)/P(B)$.\end{solution}

% ---- Q2 (false: E) ----
\question[3] (Select all that apply) Which statements follow from the standard
rules of probability?
\begin{checkboxes}
  \CorrectChoice Product rule: $P(X,Y)=P(X\mid Y)\,P(Y)$.
  \CorrectChoice Marginalization: $P(X{=}x)=\sum_{y}P(X{=}x,Y{=}y)$.
  \CorrectChoice Law of total probability: for a partition $\{B_i\}$,
        $P(A)=\sum_i P(A\mid B_i)\,P(B_i)$.
  \CorrectChoice Complement rule: $P(\lnot A)=1-P(A)$.
  \choice If $P(A\mid B)=P(A)$, then $A$ and $B$ must be mutually exclusive.
  \CorrectChoice For any event $A$, $0\le P(A)\le 1$.
\end{checkboxes}
\begin{solution}
\textbf{Correct: A, B, C, D, F.} The false option \textbf{E}: $P(A\mid B)=P(A)$
means $A$ and $B$ are \emph{independent}, not mutually exclusive.
\end{solution}

% ---- Q3 (a: B, b: C) ----
\question[5] (Calculation) In a factory, Machine~1 makes $60\%$ of the parts and
$3\%$ of its parts are defective; Machine~2 makes the other $40\%$ and $5\%$ of
its parts are defective. A part is drawn at random and found defective.
\begin{parts}
  \part Which expression gives the overall probability a part is defective,
        $P(\text{def})$?
  \begin{choices}
    \choice $0.60+0.40$
    \CorrectChoice $0.60(0.03)+0.40(0.05)$
    \choice $0.03+0.05$
    \choice $0.60(0.03)\times0.40(0.05)$
  \end{choices}
  \begin{solution}\textbf{Law of total probability}; value $=0.018+0.020=0.038$.\end{solution}

  \part Using $P(\text{def})=0.038$, which expression gives
        $P(\text{Machine 2}\mid\text{def})$, and its value?
  \begin{choices}
    \choice $\dfrac{0.038}{0.40(0.05)}\approx 1.90$
    \choice $\dfrac{0.40}{0.038}\approx 10.5$
    \CorrectChoice $\dfrac{0.40(0.05)}{0.038}\approx 0.53$
    \choice $0.40(0.05)\times0.038\approx 0.0008$
  \end{choices}
  \begin{solution}\textbf{Bayes rule}: $0.020/0.038=0.526\approx0.53$.\end{solution}
\end{parts}

%==============================================================================
\topic{2}{Bayesian Network: Representation}
%==============================================================================

% ---- Q1 (correct: A) ----
\question[2] (Single answer) A node is conditionally independent of all other
nodes in the network given which set?
\begin{choices}
  \CorrectChoice Its Markov blanket.
  \choice Its parents and its children only.
  \choice Every other node, regardless of structure.
  \choice Its ancestors together with its descendants.
\end{choices}
\begin{solution}\textbf{A.} A node's Markov blanket -- its parents, its children,
and its children's other parents -- shields it from the rest of the network.\end{solution}

% ---- Q2 (false: C) ----
\question[3] (Select all that apply) Which statements about Bayesian-network
\emph{representation} are correct?
\begin{checkboxes}
  \CorrectChoice The joint factorizes as a product of each node's conditional
        given its parents.
  \CorrectChoice A binary node with $k$ binary parents needs $2^{k}$ independent
        parameters (one per parent configuration).
  \choice Adding more edges to a network always lets it represent \emph{fewer}
        joint distributions.
  \CorrectChoice Two networks with different structures can encode the same
        independencies (they are Markov equivalent).
  \CorrectChoice A node is independent of its non-descendants given its parents.
  \CorrectChoice In general, the joint over $n$ binary variables needs $2^{n}-1$
        free parameters.
\end{checkboxes}
\begin{solution}
\textbf{Correct: A, B, D, E, F.} The false option \textbf{C}: more edges impose
\emph{fewer} independence constraints, so the network can represent \emph{more}
distributions, not fewer.
\end{solution}

% ---- Q3 (a: D, b: A) ----
\question[5] (Calculation) Consider the three-node chain below (each arrow points
from a cause to its effect):
\[
  C \longrightarrow R \longrightarrow W .
\]
In words: Cloudy $C$ influences Rain $R$, which influences a Wet lawn $W$, so the
joint factorizes as $P(C,R,W)=P(C)\,P(R\mid C)\,P(W\mid R)$. The values are
$P(C{=}1)=0.5$, $P(R{=}1\mid C{=}1)=0.8$, and $P(W{=}1\mid R{=}1)=0.9$.
\begin{parts}
  \part Which expression gives $P(C{=}1,R{=}1,W{=}1)$?
  \begin{choices}
    \choice $P(C{=}1)+P(R{=}1\mid C{=}1)+P(W{=}1\mid R{=}1)$
    \choice $P(C{=}1)\times P(R{=}1)\times P(W{=}1)$
    \choice $P(R{=}1\mid C{=}1)\times P(W{=}1\mid R{=}1)$
    \CorrectChoice $P(C{=}1)\times P(R{=}1\mid C{=}1)\times P(W{=}1\mid R{=}1)$
  \end{choices}
  \begin{solution}\textbf{Product of the three chain factors} (chain rule along the DAG).\end{solution}

  \part That product $0.5\times0.8\times0.9$ equals:
  \begin{choices}
    \CorrectChoice $0.36$
    \choice $0.72$
    \choice $0.045$
    \choice $2.2$
  \end{choices}
  \begin{solution}\textbf{0.36} $=0.5\times0.8\times0.9$.\end{solution}
\end{parts}

%==============================================================================
\topic{3}{Bayesian Network: Inference}
%==============================================================================

% ---- Q1 (correct: D) ----
\question[2] (Single answer) Variable elimination is usually far cheaper than
enumerating the full joint because it:
\begin{choices}
  \choice always runs in linear time on every possible network.
  \choice first builds the complete joint table and then deletes rows.
  \choice works only when all variables are mutually independent.
  \CorrectChoice sums out variables one at a time, reusing intermediate factors.
\end{choices}
\begin{solution}\textbf{D.} It eliminates variables sequentially and reuses factors,
never materializing the full joint.\end{solution}

% ---- Q2 (false: F) ----
\question[3] (Select all that apply) Which statements about Bayesian-network
\emph{inference} are correct?
\begin{checkboxes}
  \CorrectChoice Marginalization sums out variables that are neither queried nor
        observed.
  \CorrectChoice Normalizing (conditioning) divides by the probability of the
        evidence.
  \CorrectChoice The running time of variable elimination depends on the
        elimination ordering.
  \CorrectChoice Exact inference in general Bayesian networks is NP-hard.
  \CorrectChoice Variable elimination reuses intermediate factors rather than
        rebuilding the joint.
  \choice An optimal elimination ordering can always be found in polynomial time.
\end{checkboxes}
\begin{solution}
\textbf{Correct: A, B, C, D, E.} The false option \textbf{F}: finding an optimal
elimination ordering is itself NP-hard.
\end{solution}

% ---- Q3 (a: C, b: B) ----
\question[5] (Calculation) A message is Spam with prior $P(\text{Spam})=0.4$. A
particular word appears with $P(\text{word}\mid\text{Spam})=0.6$ and
$P(\text{word}\mid\text{not Spam})=0.2$. A message contains the word.
\begin{parts}
  \part Which expression gives $P(\text{word})$ (marginalizing over Spam / not Spam)?
  \begin{choices}
    \choice $0.6+0.2$
    \choice $0.4(0.6)\times0.6(0.2)$
    \CorrectChoice $0.4(0.6)+0.6(0.2)$
    \choice $0.4(0.6)$
  \end{choices}
  \begin{solution}\textbf{Law of total probability}; value $=0.24+0.12=0.36$.\end{solution}

  \part Using $P(\text{word})=0.36$, which expression gives
        $P(\text{Spam}\mid\text{word})$, and its value?
  \begin{choices}
    \choice $\dfrac{0.6(0.2)}{0.36}\approx 0.33$
    \CorrectChoice $\dfrac{0.4(0.6)}{0.36}\approx 0.67$
    \choice $\dfrac{0.36}{0.4(0.6)}\approx 1.5$
    \choice $\dfrac{0.4}{0.36}\approx 1.11$
  \end{choices}
  \begin{solution}\textbf{Bayes rule}: $0.24/0.36=0.667\approx0.67$.\end{solution}
\end{parts}

%==============================================================================
\topic{4}{Markov Decision Processes}
%==============================================================================

% ---- Q1 (correct: B) ----
\question[2] (Single answer) What is the role of the discount factor
$\gamma\in[0,1)$ in an MDP?
\begin{choices}
  \choice It sets the probability of taking a random action.
  \CorrectChoice It weights immediate rewards more heavily than distant future rewards.
  \choice It is the learning rate used by value iteration.
  \choice It guarantees every reward is positive.
\end{choices}
\begin{solution}\textbf{B.} $\gamma$ discounts future rewards; smaller $\gamma$ is
more short-sighted, and $\gamma<1$ also keeps the infinite-horizon return finite.\end{solution}

% ---- Q2 (false: B) ----
\question[3] (Select all that apply) Which statements about MDPs are correct?
\begin{checkboxes}
  \CorrectChoice An optimal policy $\pi^{*}$ picks, in each state, an action that
        maximizes expected return.
  \choice A \emph{larger} discount factor $\gamma$ makes the agent more
        short-sighted, ignoring future rewards.
  \CorrectChoice The Bellman optimality equation relates a state's value to the
        values of its successor states.
  \CorrectChoice Value iteration repeatedly applies the Bellman update until the
        values converge.
  \CorrectChoice Policy iteration alternates policy evaluation with policy
        improvement.
  \CorrectChoice $V^{*}(s)=\max_a Q^{*}(s,a)$.
\end{checkboxes}
\begin{solution}
\textbf{Correct: A, C, D, E, F.} The false option \textbf{B} is backwards: a larger
$\gamma$ makes the agent more \emph{far-sighted}.
\end{solution}

% ---- Q3 (a: A, b: D) ----
\question[5] (Calculation) From state $s$ an agent may take Action~1 or Action~2,
with discount $\gamma=0.5$. Action~1 gives reward $3$ and leads to a state of
value $V=10$. Action~2 gives reward $6$ and leads to a state of value $V=2$.
\begin{parts}
  \part Which expression gives the value $Q(s,\text{Action 1})$?
  \begin{choices}
    \CorrectChoice $3+0.5(10)$
    \choice $3\times0.5\times10$
    \choice $3+10$
    \choice $0.5(3)+10$
  \end{choices}
  \begin{solution}\textbf{$Q=r+\gamma V(s')$} $=3+0.5(10)=8$.\end{solution}

  \part Since $Q(s,\text{Action 1})=8$ and $Q(s,\text{Action 2})=6+0.5(2)=7$, the
        state value $V(s)=\max_a Q(s,a)$ and the chosen action are:
  \begin{choices}
    \choice $7$; choose Action 2.
    \choice $6.5$; either action.
    \choice $15$; choose Action 1.
    \CorrectChoice $8$; choose Action 1.
  \end{choices}
  \begin{solution}\textbf{$V(s)=8$, Action 1} -- the maximum-value action ($8>7$).\end{solution}
\end{parts}

%==============================================================================
\topic{5}{Reinforcement Learning}
%==============================================================================

% ---- Q1 (correct: C) ----
\question[2] (Single answer) What distinguishes an \emph{off-policy} method such as
Q-learning from an \emph{on-policy} method such as Sarsa?
\begin{choices}
  \choice Off-policy methods require a full model of the environment's dynamics.
  \choice On-policy methods never explore.
  \CorrectChoice Off-policy methods evaluate one policy while following another.
  \choice On-policy methods do not use a value function.
\end{choices}
\begin{solution}\textbf{C.} Q-learning learns about the greedy (target) policy while
acting under a different behaviour policy; Sarsa learns about the policy it follows.\end{solution}

% ---- Q2 (false: D) ----
\question[3] (Select all that apply) Which statements about reinforcement learning
are correct?
\begin{checkboxes}
  \CorrectChoice RL learns from sampled experience when the transition and reward
        models are unknown.
  \CorrectChoice Temporal-difference methods update after each step using a
        bootstrapped estimate.
  \CorrectChoice Monte Carlo methods wait until an episode ends before updating.
  \choice Raising $\epsilon$ in $\epsilon$-greedy always increases the final
        policy's reward.
  \CorrectChoice Q-learning is off-policy, whereas Sarsa is on-policy.
  \CorrectChoice The exploration--exploitation trade-off governs how often the
        agent tries new actions.
\end{checkboxes}
\begin{solution}
\textbf{Correct: A, B, C, E, F.} The false option \textbf{D}: too much exploration
(large $\epsilon$) can \emph{lower} performance.
\end{solution}

% ---- Q3 (a: C, b: B) ----
\question[5] (Calculation) An agent using \textbf{Sarsa} observes reward $r=3$,
then in the next state takes the action whose value is $Q(s',a')=8$. Currently
$Q(s,a)=4$, with learning rate $\alpha=0.5$ and discount $\gamma=0.5$.
\begin{parts}
  \part Which expression is the correct Sarsa update for $Q(s,a)$?
  \begin{choices}
    \choice $Q(s,a)+\alpha\bigl[r+\gamma\max_{a'}Q(s',a')-Q(s,a)\bigr]$
    \choice $Q(s,a)+\alpha\bigl[r+\gamma\,Q(s',a')\bigr]$
    \CorrectChoice $Q(s,a)+\alpha\bigl[r+\gamma\,Q(s',a')-Q(s,a)\bigr]$
    \choice $r+\gamma\,Q(s',a')$
  \end{choices}
  \begin{solution}\textbf{C} -- Sarsa uses the value of the action actually taken
  next, $Q(s',a')$, not the $\max$ (that would be Q-learning).\end{solution}

  \part Substituting, $4+0.5\bigl[3+0.5(8)-4\bigr]$ equals:
  \begin{choices}
    \choice $7.0$
    \CorrectChoice $5.5$
    \choice $3.5$
    \choice $4.0$
  \end{choices}
  \begin{solution}\textbf{5.5} (target $=3+4=7$, so $4+0.5(7-4)=5.5$).\end{solution}
\end{parts}

%==============================================================================
\topic{6}{Heuristic Search}
%==============================================================================

% ---- Q1 (correct: A) ----
\question[2] (Single answer) In A$^{*}$ search with $f(n)=g(n)+h(n)$, what do
$g(n)$ and $h(n)$ denote?
\begin{choices}
  \CorrectChoice $g$ = cost so far (start $\to n$); $h$ = estimated cost to go ($n\to$ goal).
  \choice $g$ = estimated cost to the goal; $h$ = the exact cost already paid from the start.
  \choice Both are always equal to zero.
  \choice $g$ = the goal depth; $h$ = the branching factor.
\end{choices}
\begin{solution}\textbf{A.} $g$ is the known cost from the start; $h$ estimates the
remaining cost to the goal.\end{solution}

% ---- Q2 (false: A) ----
\question[3] (Select all that apply) Which statements about search algorithms are
correct?
\begin{checkboxes}
  \choice Greedy best-first search always returns the least-cost path to the goal.
  \CorrectChoice A$^{*}$ is optimal when the heuristic is admissible (and, for
        graph search, consistent).
  \CorrectChoice Breadth-first search is complete and finds the shallowest goal.
  \CorrectChoice A consistent heuristic is also admissible.
  \CorrectChoice Uniform-cost search expands nodes in order of path cost $g(n)$.
  \CorrectChoice Setting $h(n)=0$ for all $n$ turns A$^{*}$ into uniform-cost search.
\end{checkboxes}
\begin{solution}
\textbf{Correct: B, C, D, E, F.} The false option \textbf{A}: greedy best-first
search is not optimal -- it follows $h$ alone and can miss the cheapest path.
\end{solution}

% ---- Q3 (a: B, b: D) ----
\question[5] (Calculation) During an A$^{*}$ search, three nodes sit on the
frontier with these cost-so-far $g$ and heuristic $h$ values:
node~$A$ has $g=2,\ h=6$; node~$B$ has $g=4,\ h=3$; node~$C$ has $g=5,\ h=5$.
\begin{parts}
  \part Which expression gives $f(B)$?
  \begin{choices}
    \choice $4\times3$
    \CorrectChoice $4+3$
    \choice $\max(4,3)$
    \choice $4+3+5$
  \end{choices}
  \begin{solution}\textbf{$f(B)=g(B)+h(B)=4+3=7$}.\end{solution}

  \part A$^{*}$ expands the frontier node with the smallest $f$. Given
        $f(A)=8$, $f(B)=7$, $f(C)=10$, the node expanded next is:
  \begin{choices}
    \choice $A$, with $f=8$.
    \choice $C$, with $f=10$.
    \choice $A$, because it has the smallest $g$.
    \CorrectChoice $B$, with $f=7$.
  \end{choices}
  \begin{solution}\textbf{$B$ ($f=7$)} -- the smallest $f$-value on the frontier.\end{solution}
\end{parts}

%==============================================================================
\topic{7}{Deep Q-learning}
%==============================================================================

% ---- Q1 (correct: D) ----
\question[2] (Single answer) Why does DQN use a separate \textbf{target network}?
\begin{choices}
  \choice To eliminate the replay buffer and store nothing between updates.
  \choice To make the network smaller and run faster.
  \choice To turn the task into ordinary i.i.d.\ supervised learning.
  \CorrectChoice To keep the update targets stable across steps.
\end{choices}
\begin{solution}\textbf{D.} Freezing the target network for several steps stops the
target from shifting every update, which stabilizes training.\end{solution}

% ---- Q2 (false: E) ----
\question[3] (Select all that apply) Which statements about Deep Q-learning are
correct?
\begin{checkboxes}
  \CorrectChoice A replay buffer stores past transitions and samples mini-batches
        to decorrelate updates.
  \CorrectChoice A target network holds the bootstrapped targets fixed for several
        steps to stabilize learning.
  \CorrectChoice DQN can learn a control policy directly from high-dimensional
        pixel input.
  \CorrectChoice Double DQN reduces the overestimation bias of the $\max$ operator.
  \choice Because the targets are held fixed, DQN training is exactly ordinary
        i.i.d.\ supervised regression.
  \CorrectChoice The loss is a squared temporal-difference error between the target
        and the current $Q$-value.
\end{checkboxes}
\begin{solution}
\textbf{Correct: A, B, C, D, F.} The false option \textbf{E}: the data are
correlated and non-stationary, unlike an i.i.d.\ supervised dataset.
\end{solution}

% ---- Q3 (a: A, b: C) ----
\question[5] (Calculation) A DQN transition has reward $r=2$ and discount
$\gamma=0.9$. For the next state, the target network's best value is
$\max_{a'}\hat Q(s',a')=10$.
\begin{parts}
  \part Which expression gives the temporal-difference \emph{target} $y$?
  \begin{choices}
    \CorrectChoice $r+\gamma\,\max_{a'}\hat Q(s',a')$
    \choice $r+\gamma\,\hat Q(s',a)-Q(s,a)$
    \choice $r\times\gamma\times\max_{a'}\hat Q(s',a')$
    \choice $\max_{a'}\hat Q(s',a')-r$
  \end{choices}
  \begin{solution}\textbf{A} -- the DQN target bootstraps from the target network's
  best next-state value.\end{solution}

  \part Substituting, $2+0.9(10)$ equals:
  \begin{choices}
    \choice $9.0$
    \choice $20$
    \CorrectChoice $11$
    \choice $10.8$
  \end{choices}
  \begin{solution}\textbf{11} $=2+9$.\end{solution}
\end{parts}

%==============================================================================
\topic{8}{Policy Learning}
%==============================================================================

% ---- Q1 (correct: B) ----
\question[2] (Single answer) Why is a \textbf{baseline} subtracted in policy-gradient
methods such as REINFORCE?
\begin{choices}
  \choice To make the learned policy fully deterministic over time.
  \CorrectChoice To lower the gradient's variance while keeping it unbiased.
  \choice To increase the learning rate automatically.
  \choice To remove the need for any reward signal.
\end{choices}
\begin{solution}\textbf{B.} A baseline that does not depend on the action reduces
variance without changing the gradient's expectation.\end{solution}

% ---- Q2 (false: C) ----
\question[3] (Select all that apply) Which statements about policy-gradient methods
are correct?
\begin{checkboxes}
  \CorrectChoice The policy $\pi_\theta(a\mid s)$ must be differentiable with
        respect to $\theta$.
  \CorrectChoice A state-dependent baseline reduces variance without biasing the
        gradient.
  \choice A baseline that depends on the chosen \emph{action} still leaves the
        gradient unbiased.
  \CorrectChoice Actor-critic methods learn a value estimate (critic) alongside the
        policy (actor).
  \CorrectChoice The advantage $A(s,a)=Q(s,a)-V(s)$ is positive for
        better-than-average actions.
  \CorrectChoice REINFORCE is a Monte Carlo policy-gradient method.
\end{checkboxes}
\begin{solution}
\textbf{Correct: A, B, D, E, F.} The false option \textbf{C}: the baseline must
\emph{not} depend on the action, or it biases the gradient.
\end{solution}

% ---- Q3 (a: D, b: A) ----
\question[5] (Calculation) A plain \textbf{REINFORCE} update (no baseline) uses
$\theta_t=1.0$, learning rate $\alpha=0.2$, return $G_t=5$, and
$\nabla_\theta\ln\pi_\theta(A_t\mid S_t)=3$.
\begin{parts}
  \part Which expression is the correct REINFORCE update for $\theta_{t+1}$?
  \begin{choices}
    \choice $\theta_t-\alpha\,G_t\,\nabla_\theta\ln\pi_\theta(A_t\mid S_t)$
    \choice $\theta_t+\alpha\,\nabla_\theta\ln\pi_\theta(A_t\mid S_t)$
    \choice $\theta_t+\alpha\,G_t$
    \CorrectChoice $\theta_t+\alpha\,G_t\,\nabla_\theta\ln\pi_\theta(A_t\mid S_t)$
  \end{choices}
  \begin{solution}\textbf{D} -- ascend the gradient, scaled by the return $G_t$.\end{solution}

  \part Substituting, $1.0+0.2(5)(3)$ equals:
  \begin{choices}
    \CorrectChoice $4.0$
    \choice $3.0$
    \choice $1.6$
    \choice $7.0$
  \end{choices}
  \begin{solution}\textbf{4.0} $=1.0+0.2\times15=1.0+3.0$.\end{solution}
\end{parts}

%==============================================================================
\topic{9}{RL for Large Language Models}
%==============================================================================

% ---- Q1 (correct: A) ----
\question[2] (Single answer) In reinforcement learning from human feedback (RLHF),
what is the purpose of the \textbf{reward model}?
\begin{choices}
  \CorrectChoice To score how well a response matches human preferences.
  \choice To store the entire pre-training corpus for later reuse.
  \choice To translate text between different natural languages.
  \choice To compress the policy network into a smaller lookup table.
\end{choices}
\begin{solution}\textbf{A.} The reward model predicts a human-preference score, which
then guides the policy's optimization.\end{solution}

% ---- Q2 (false: D) ----
\question[3] (Select all that apply) The statements below concern aligning a large
language model with human feedback. (Reminder: \emph{KL-divergence} measures how
different two probability distributions are.) Which statements are correct?
\begin{checkboxes}
  \CorrectChoice Supervised fine-tuning teaches a pre-trained model to follow
        instructions.
  \CorrectChoice A reward model is trained from human preference comparisons and
        then scores candidate responses.
  \CorrectChoice Proximal Policy Optimization (PPO) uses a KL-divergence penalty to
        keep the updated policy close to the reference model.
  \choice Direct Preference Optimization (DPO) still requires training and querying
        a separate reward model at every step.
  \CorrectChoice Group Relative Policy Optimization (GRPO) replaces the learned value
        function with a baseline from a group of sampled responses.
  \CorrectChoice The goal of alignment is to make model outputs better match human
        preferences.
\end{checkboxes}
\begin{solution}
\textbf{Correct: A, B, C, E, F.} The false option \textbf{D}: DPO optimizes the
preference objective \emph{directly}, with no separate reward model.
\end{solution}

% ---- Q3 (a: B, b: C) ----
\question[5] (Calculation) During RLHF, a response's raw reward is combined with a
penalty that discourages drifting from the reference model. Here the raw reward is
$r=5$, the penalty weight is $\beta=0.2$, and the measured KL-divergence from the
reference is $\mathrm{KL}=4$.
\begin{parts}
  \part Which expression gives the KL-penalized reward?
  \begin{choices}
    \choice $r+\beta\cdot\mathrm{KL}$
    \CorrectChoice $r-\beta\cdot\mathrm{KL}$
    \choice $\beta\cdot\mathrm{KL}-r$
    \choice $r\times\beta\times\mathrm{KL}$
  \end{choices}
  \begin{solution}\textbf{$r-\beta\,\mathrm{KL}$} -- the penalty is \emph{subtracted}
  to discourage large drift from the reference model.\end{solution}

  \part Substituting, $5-0.2(4)$ equals:
  \begin{choices}
    \choice $5.8$
    \choice $4.0$
    \CorrectChoice $4.2$
    \choice $0.8$
  \end{choices}
  \begin{solution}\textbf{4.2} $=5-0.8$.\end{solution}
\end{parts}

\end{questions}

\vfill
\begin{center}\small\itshape End of assessment --- 9 topics, 3 questions each; all parts multiple choice.\end{center}

\end{document}
